\clearpage
\fancyhead[L]{\small STEINHEIL. PRISMENPHOTOMETER.}
\SourceLeaf{654}
\begin{center}\small Göttingische gelehrte Anzeigen. Stück 34. 35. Seiten 329..337. 5. März 1835.\\[3pt]\rule{.72\linewidth}{.3pt}\end{center}

Göttingen. Zur Beantwortung der für November 1834 von der Mathematischen Klasse der Königlichen Sozietät aufgegebene Hauptpreisfrage, deren Frist aber laut Stück 149 dieser Anzeigen vom Vorjahr bis Ende Dezember verlängert worden war, waren drei Konkurrenzschriften eingegangen. Eine war in lateinischer Sprache mit dem Motto: Opinionum commenta delet dies, naturae iudicia confirmat. Die zweite war in deutscher Sprache mit der Aufschrift: Suum cuique. Die dritte war ebenfalls auf Deutsch mit den Worten: Nur gleichartige Eindrücke sind vergleichbar.

Die Abhandlung Nr.~2 mit der Aufschrift Suum cuique enthält nur die Meinungen ihres Verfassers über die Entstehung und natürliche Beschaffenheit der Himmelskörper und gar nichts, was sich auf die Lösung der von der Sozietät gestellten Aufgabe bezieht. Eine besondere Beurteilung dieser Meinungen ist daher unnötig, da sie mit der Preisfrage in gar keinem Zusammenhang stehen.

Der Verfasser der Schrift Nr.~1, Opinionum commenta usw., hat dagegen die Frage richtig aufgefasst, einen Apparat zum Vergleich der Lichtstärke zweier Sterne angegeben und bauen lassen sowie einige Versuche zur Anwendung auf wirkliche Lichtmessungen mitgeteilt. Das Instrument ist ein Fernrohr mit Vorrichtungen, durch die beide Sterne zugleich im Gesichtsfeld nebeneinander gesehen werden können, der eine direkt, der andere durch Reflexion. Diese wird durch einen vor dem Objektiv angebrachten Spiegel bewirkt. Durch Drehung um eine Achse, die die Gesichtslinie rechtwinklicht schneidet, lässt sich der Spiegel in die dem Winkelabstand beider Sterne entsprechende Neigung zur Gesichtslinie bringen. Der äußere Rand des Spiegels fällt mit dieser Drehachse zusammen, weshalb der Spiegel in jeder Lage die Hälfte des Spiegels für direktes Licht verschattet. Unmittelbar vor dem Objektiv ist außerdem eine halbkreisförmige Blende angebracht, die nur die Hälfte des Objektivs offen lässt und ganz herumgedreht werden kann. Die Größe dieser Drehung wird an einem eingeteilten Ring gemessen (wie die Größe der Spiegeldrehung an einem Gradbogen). Steht der Index des Rings auf dem Nullpunkt, gelangt gar kein direktes Licht ins Fernrohr, nach einer halben Umdrehung dagegen gar kein reflektiertes Licht. In jeder Zwischenstellung teilen sich das reflektierte und das direkte Licht die offene Hälfte des Objektivs im Verhältnis der Abweichungen von jenen beiden Stellungen. Man erkennt also leicht, dass sich, wenn man durch Drehen der Objektivblende erreicht hat, dass beide Sterne gleich hell erscheinen, das Verhältnis der Lichtstärken beider Sterne bis auf einen noch unbekannten Faktor berechnen lässt, der von der Abschwächung des Lichts durch die Reflexion abhängt. Dieser unbekannte Faktor wird durch Heranziehen einer zweiten Beobachtung gefunden oder eliminiert, bei der bloß die Sterne vertauscht werden. Für bestimmte Fälle hat der Verfasser noch einen zweiten Spiegel hinzugefügt, sodass der eine Stern durch doppelte Reflexion gesehen wird, was im Übrigen für die Methode keinen Unterschied macht. Ein parallaktisches Stativ macht den Gebrauch erheblich bequemer.

Es ist zu bedauern, dass der Verfasser dieses Instrument so spät vom Hersteller erhielt, dass er keine gründliche Prüfung durch zahlreiche Messungen vornehmen konnte. Er hat das Lichtverhältnis von sieben Sternpaaren aus insgesamt nur 44 Beobachtungen bestimmt, die jedoch nur zusammenfassend angegeben werden. Die Ergebnisse, bei denen jeweils der zuerst genannte Stern als Einheit betrachtet wird, sind folgende.
\SourceLeaf{655}
\begin{center}
\begin{tabular}{lc}
Sterne & Lichtverhältnis\\[3pt]
Rigel, Procyon & $0.8501$\\
Rigel, $\beta$ kl. Hund & $0.1258$\\
Sirius. Rigel & $0.2875$\\
Sirius, Procyon & $0.2756$\\
Procyon, Regulus & $0.3781$\\
Procyon, Nordstern & $0.4369$\\
Regulus, Nordstern & $0.5720$
\end{tabular}
\end{center}

Die Höhen der Sterne oder die Größen, von denen sie abhängen, fehlen. Die wahrscheinlichen Fehler dieser Bestimmungen würden nach den Angaben des Verfassers, soweit sie durch den Vergleich der einzelnen Beobachtungen untereinander bestimmt werden können, zwischen $\frac{1}{16}$ und $\frac{1}{64}$ des Ganzen schwanken. Vergleicht man aber die erste, dritte und vierte Bestimmung miteinander, zeigt sich, dass viel stärkere Korrekturen nötig sind, und die drei letzten Bestimmungen lassen sich überhaupt nicht miteinander vereinbaren. Der Verfasser räumt selbst ein, dass er diesen Widerspruch nicht zu erklären weiß. Auch wenn man hoffen muss, dass es ihm in Zukunft nach viel umfassenderen Versuchen gelingen wird, die Quelle solcher Fehler zu finden, bleibt die Eignung des Apparats zur Messung der Helligkeit leuchtender Punkte gegenwärtig noch unverbürgt.

Der Verfasser der dritten Abhandlung mit dem Motto \emph{Nur gleichartige Eindrücke sind vergleichbar} hat zwei ganz verschiedene Apparate angegeben und gebaut. Den einen nennt er den Okularapparat, den anderen das Prismenphotometer. Obwohl beide für den vorgegebenen Zweck verwendet werden können, ist der erste eigentlich weniger zum Vergleich der Lichtstärke leuchtender Punkte als zum Vergleich der spezifischen Helligkeit ausgedehnterer Flächen, zum Beispiel des Himmelsgrundes, bestimmt. Es genügt daher, hier nur die wesentlichen Merkmale des zweiten Apparats anzugeben. Der Grundgedanke dieses Instruments ist die bekannte Erfahrung, dass ein Stern, der dem unbewaffneten Auge oder in einem zum deutlichen Sehen eingestellten Fernrohr als unteilbarer leuchtender Punkt erscheint, sich in ein kreisförmigs Bild ausbreitet, wenn man dem Okular eine andere Stellung gibt, als das deutliche Sehen erfordert. Je weiter das Okular von seiner Normalstellung entfernt ist, desto größer, aber gerade deshalb in seinen Teilen desto lichtschwächer ist dieses Bild. Bei unterschiedlich hellen Sternen muss man das Okular daher in unterschiedliche Entfernungen von der Normalstellung bringen, damit die Bilder mit gleicher Flächenhelligkeit erscheinen. So lässt sich die Lichtstärke zweier Sterne schon einigermaßen vergleichen, wenn man unscharfe Bilder von ihnen \emph{nacheinander} beobachtet, ihre Flächenhelligkeit, soweit es der im Gedächtnis gebliebene Eindruck erlaubt, gleich macht und die entsprechenden Okularstellungen misst. Natürlich erwartet man von einem so groben Verfahren wenig Genauigkeit und ist daher überrascht, dass die vom Verfasser angeführten Versuche eine viel größere Übereinstimmung zeigen, als man hätte erwarten können. Das weckt schon ein günstiges Vorurteil für den vom Verfasser kunstreich angeordneten Apparat, mit dem man solche Bilder zweier Sterne \emph{zugleich} sehen und auf gleiche Flächenhelligkeit bringen kann.

Das Objektiv ist in zwei gleiche Hälften zerschnitten. Anders als beim Heliometer lassen sich die Hälften einzeln längs ihrer gemeinsamen Achse statt nebeneinander verschieben. Die Mitte der Verschiebungen, die durch Skalen an der Außenseite des Rohrs genau gemessen werden, entspricht bei ganz eingeschobener Okularröhre ungefähr der Stellung gegenüber dem Rohr, die zum deutlichen Sehen erforderlich ist. Die beiden Objektivhälften erhalten ihr Licht durch Spiegel, deren reflektierende Flächen unter einem Winkel von $45^\circ$ zur Achse des Rohrs geneigt sind. Einer der Spiegel (der weiter vom Objektiv entfernte) kann um diese Achse messbar gedreht werden. Beim Beobachten zweier Sterne ist diese Achse also immer auf einen Pol des Großkreises zu richten, der sie verbindet. Die \SourceLeaf{656}Spiegel selbst sind Glasprismen, in die das Licht senkrecht einfällt und aus denen es senkrecht austritt. Zwischen den Objektivhälften und den zugehörigen Prismenspiegeln sind Diaphragmen angebracht, die durch zwei Schieberpaare gebildet werden. Jedes Schieberpaar wird durch eine einzige Schraube mit entgegengesetzt geschnittenen Gewinden so bewegt, dass die Mitte der Hypotenuse des Diaphragmas unverrückt bleibt, wenn dieses ein größeres oder kleineres rechtwinkliges Dreieck bildet.

Mit dieser Anordnung sieht man bei geeigneter Stellung des Rohrs und der Spiegel zwei Sterne zugleich, und zwar jeden als rechtwinklige Dreiecksfläche, wenn die Objektivhälften von ihrer Normallage zum Okular abweichen. Von dieser Abweichung hängen sowohl die scheinbare Größe des Dreiecks als auch seine Flächenhelligkeit ab. Die Größe hängt außerdem von der Diaphragmenöffnung, die Flächenhelligkeit von der eigenen Helligkeit jedes Sterns ab. Man kann daher durch Änderung einer der beiden Abweichungen die Flächenhelligkeiten beider Bilder und, wenn man will, durch Änderung einer Diaphragmenöffnung auch ihre Größe einander angleichen. Dass sich so das Verhältnis der Lichtstärke zweier Sterne bestimmen lässt und dabei etwaige Ungleichheiten in den Objektivhälften und Prismenspiegeln durch Umkehrung der Kombination beseitigt werden können, bedarf keiner weiteren Ausführung.

Der Verfasser hat seinen Apparat einer strengen Prüfung unterzogen, und zwar bewusst an künstlich erzeugten sternähnlichen Leuchtpunkten statt an Sternen. Diese künstlichen Sterne erhielt er durch die Reflexion des Tageslichts an zwei nahezu gleichen, gut polierten Stahlkugeln von etwa $\frac{1}{2}$ Zoll Durchmesser. Das Tageslicht kam für beide Kugeln von derselben Stelle des Himmelshintergrunds und gelangte durch kreisrunde Blenden unterschiedlicher Weite zu den Kugeln. Es wurde dafür gesorgt, dass kein Fremdlicht die Kugeln oder das Auge des Beobachters treffen konnte. Insgesamt wurden vier Blenden verwendet, die engste mit 7, die weiteste mit 20 Linien Durchmesser. Durch die sechs verschiedenen Kombinationen konnte man also künstliche Sterne mit sechs verschiedenen Lichtverhältnissen erhalten. Die größte Ungleichheit, 1 zu 8, entspricht nach den eigenen Untersuchungen des Verfassers annähernd dem mittleren Lichtverhältnis zweier Sterne, die zwei Größenklassen auseinanderliegen. Diese künstlichen Sterne erscheinen wirklich ganz ähnlich, jedoch ohne die Schwankungen und das Wallen, durch die Beobachtungen wirklicher Sterne oft so unsicher werden. Außerdem hatten sie den äußerst wichtigen Vorzug, dass ihr Helligkeitsverhältnis aus den Blendenöffnungen a priori bekannt war. Der Verfasser teilt die zahlreichen Messungen ihrer Lichtstärke mit dem Prismenphotometer ausführlich mit, ohne diejenigen zu verschweigen, bei denen sich anfangs einige Unregelmäßigkeiten zeigten, deren Ursachen jedoch entdeckt und beseitigt wurden. Aus der Gesamtheit der Messungen ergibt sich der wahrscheinliche Fehler eines einzigen Vergleichs als $\frac{1}{32}$ der ganzen Helligkeit, ob diese nun groß oder klein ist. Die Verhältnisse der verschiedenen künstlichen Sterne zeigen eine vollkommen befriedigende Übereinstimmung mit den Blendenöffnungen.

Die Eignung des Apparats zum genauen Vergleich der Helligkeit leuchtender Punkte ist damit hinreichend erwiesen. Auch wenn man das Fehlen von Anwendungen auf wirkliche Sterne bedauert, besteht doch Grund genug, auch bei ihnen befriedigende Ergebnisse zu erwarten, sofern man, wie der Verfasser zu Recht verlangt, die Beobachtungen auf besonders günstige atmosphärische Bedingungen beschränkt. Unter solchen Bedingungen wird man dank der leichten Handhabung des Instruments in wenigen Stunden mehr ausrichten als unter ungünstigen Umständen an vielen Tagen. Im Übrigen enthält die Abhandlung noch manche anderen photometrischen Untersuchungen und Ansichten von erheblichem Interesse, die hier stillschweigend übergangen werden können, da sie für die Hauptsache nicht wesentlich sind. Einige Anwendungen der Wahrscheinlichkeitsrechnung im letzten Abschnitt bedürften einer Berichtigung, was jedoch für den Hauptgegenstand selbst ganz unwesentlich ist.

Schließlich ist noch zu bemerken, dass das Prismenphotometer, obwohl es auf einem ganz anderen
\SourceLeaf{657}Prinzip beruht als dem, das der Abhandlung Nr.~1 zugrunde liegt, zugleich die Möglichkeit bietet, Sterne nach dem anderen Prinzip zu vergleichen, nämlich durch gleichzeitig erscheinende \emph{deutliche} Bilder bei messbar verengter Objektivöffnung. Selbst bei dieser Beobachtungsart, die der Verfasser übrigens aufgrund seiner Erfahrungen für verwerflich hält, hätte die Anordnung des Prismenphotometers Vorzüge gegenüber der bei Abhandlung Nr.~1 beschriebenen.

Da die Abhandlung Nr.~3 die Aufgabe am vollkommensten und auf eine Weise gelöst hat, die einen wertvollen Fortschritt in diesem Teil der praktischen Astronomie begründet, hat die Königliche Sozietät ihr den Preis und der ebenfalls sehr verdienstvollen Abhandlung Nr.~1 das Accessit zuerkannt.

Der Verfasser der preisgekrönten Abhandlung ist laut dem in der öffentlichen Sitzung der Sozietät vom 14. Februar entsiegelten Zettel

\begin{center}
Dr.\ \textsc{Steinheil} in München.
\end{center}

Der Zettel zur Abhandlung Nr.~2 wurde in derselben Sitzung ungeöffnet verbrannt.
